\documentclass{article}
\usepackage{graphicx}
\usepackage{amsmath}
\usepackage{multirow, multicol}
\usepackage{float}

\title{CSE200: Online-1 on \LaTeX}
\author{Arunavo Anirban Aikko (2205026)}
\date{January 9, 2026}

\begin{document}

\maketitle

\section{Introduction}
Effective technical writing demands \textbf{precision} and \textit{organization} to communicate ideas clearly. In this document, we illustrate how to combine \textbf{text},
\textit{tables}, \underline{mathematical expressions}, and {\large figures} to convey information in a
structured and readable way.\\
Using \textbf{well-defined headings}, \textit{organized lists}, and \underline{emphasized keywords}
helps readers navigate content efficiently. In addition, features such as font
adjustments, inline notation ($a_i$, $b_j$), and thoughtful spacing enhance the {\large professional }appearance and consistency of the document.\\

\textbf{Overall}, a properly formatted technical report integrates clear sections,
accurate equations, and visual aids to maximize both \underline{readability} and presentation quality.

\section{Key Concepts}
\subsection{Nested and Mixed Lists Example}
\begin{enumerate}
    \item System Setup
    \begin{itemize}
        \item Install software packages
        \item Run simulation
        \item[!] Verify connections
    \end{itemize}
    \item Simulation Steps
    \begin{enumerate}
        \item[(a)] Initialize parameters
\clearpage
        \item[(b)] Run simulation
        \begin{itemize}
            \item Monitor CPU usage
            \item Record intermediate results
        \end{itemize}
        \item[(c)] Analyze outputs 
    \end{enumerate}
    \item Reporting
    \begin{itemize}
        \item Prepare charts
        \item Document observations
        \item  Summarize findings
        \begin{enumerate}
            \item[(a)] Key Insights
            \item[(b)] Recommendations
        \end{enumerate}
    \end{itemize}
\end{enumerate}

\section{Mathematical Expressions}
Mathematical representation is crucial. Consider these examples:
\begin{equation}
    F=k\frac{q_1 q_2}{r^2}+\frac{1}{2}mv^2
\end{equation}
\begin{equation}
    \begin{split}
        g(x) &= e^x\cos(x) + ln(1+x) \\
             &= \sqrt{e^{2x}+ln^2(1+x)}
    \end{split}
\end{equation}
\begin{equation*}
    J=\int_{0}^{1}x^3e^{-x}dx
\end{equation*}
Piecewise function:
\begin{equation}
    p(t) = \begin{cases}
        \sin(t) & \text{if } 0 \le t \le \pi \\
        \ln(t+1) & \text{if } t \ge \pi
    \end{cases}
\end{equation}
\begin{equation}
    \begin{split}
        z(x) &= \sum_{n=0}^{\infty} \frac{(-1)^nx^{2n}}{(2n)!}+\frac{1}{x+1}\\
             &\quad - \sqrt[3]{1+x^2}+\frac{\tan^{-1}(x)}{2}
    \end{split}
\end{equation}
\clearpage
\section{Experiment Summary Table}
\begin{table}[H]
    \centering
    \vspace{0.2cm}
    \begin{tabular}{|l|c|c|c|c|}
        \hline
        \multirow{2}{*}{\textbf{Sample}} & \multicolumn{2}{|c|}{\textbf{Measurements}} & \multirow{2}{*}{\textbf{Status}} & \multirow{2}{*}{\textbf{Comments}} \\
        \cline{2-3} 
         & Voltage(V) & Current(A) & & \\ 
        \hline
        S1 & 5.2 & 0.8 & Pass & Stable \\ 
        \hline
        S2 & 4.9 & 0.7 & Pass & Slight fluctuation \\ 
        \hline
        \multirow{2}{*}{S3} & 3.8 & 0.6 & Fail & Low reading \\
         & 3.9 & 0.65 & Re-test & Needs adjustment \\ \hline
        S4 & 5.5 & 0.85 & Pass & Good consistency \\ \hline
    \end{tabular}
    \caption{Experimental Measurements with Multirow and Multicolumn Layout}
\end{table}
\end{document}